%%=============================================================================
%% Methodologie
%%=============================================================================

\chapter{\IfLanguageName{dutch}{Methodologie}{Methodology}}%
\label{ch:methodologie}

Dit onderzoek is opgedeeld in drie grote fasen, beginnend met de stand van zaken.

\paragraph{Fase 1}
In deze fase wordt een inleiding gemaakt voor Git en branching, met focus op deploymentprocedures, zo worden bestaande modellen en best practices onderzocht. Ook wordt er gekeken naar de huidige situatie over de mainframes van Securex. Dit resulteert in kennis van de huidige opstelling en situatie over Securex en Subversion.

\paragraph{Fase 2}
In de tweede fase wordt op basis van de bevindingen uit de literatuurstudie en context van Securex een Git-branchingstrategie ontworpen die aansluit bij de mainframe-omgeving. Hier worden er ook deploymentprocedures uitgewerkt die rekening houden met de vereisten van Securex en de doorstroom van de code. Hier wordt er ook een praktische proof-of-concept gemaakt en klein-schalig uitgevoerd in een testopstelling.

\paragraph{Fase 3}
In de derde fase wordt de ontworpen aanpak geëvalueerd op basis van vooraf gedefinieerde criteria, zoals reproduceerbaarheid, traceerbaarheid, beheersbaarheid en impact op het deploymentproces. De resultaten van de proof of concept of testopstelling worden geanalyseerd en vergeleken met de oorspronkelijke situatie.\\

Op basis van deze evaluatie worden conclusies geformuleerd en aanbevelingen gedaan voor verdere optimalisatie en mogelijke implementatie binnen Securex.

%% TODO: In dit hoofstuk geef je een korte toelichting over hoe je te werk bent
%% gegaan. Verdeel je onderzoek in grote fasen, en licht in elke fase toe wat
%% de doelstelling was, welke deliverables daar uit gekomen zijn, en welke
%% onderzoeksmethoden je daarbij toegepast hebt. Verantwoord waarom je
%% op deze manier te werk gegaan bent.
%% 
%% Voorbeelden van zulke fasen zijn: literatuurstudie, opstellen van een
%% requirements-analyse, opstellen long-list (bij vergelijkende studie),
%% selectie van geschikte tools (bij vergelijkende studie, "short-list"),
%% opzetten testopstelling/PoC, uitvoeren testen en verzamelen
%% van resultaten, analyse van resultaten, ...
%%
%% !!!!! LET OP !!!!!
%%
%% Het is uitdrukkelijk NIET de bedoeling dat je het grootste deel van de corpus
%% van je bachelorproef in dit hoofstuk verwerkt! Dit hoofdstuk is eerder een
%% kort overzicht van je plan van aanpak.
%%
%% Maak voor elke fase (behalve het literatuuronderzoek) een NIEUW HOOFDSTUK aan
%% en geef het een gepaste titel.

%\lipsum[21-25]

